\chapter{General Equilibrium with Production}

\section{Production Economy Setup}

Consider an economy with:
\begin{itemize}
    \item Two goods: $X, Y$
    \item Two factors: Labor $L$, Capital $K$
\end{itemize}

Production functions:
\[
x = F^x(L^x, K^x), \quad
y = F^y(L^y, K^y)
\]

Resource constraints:
\[
L^x + L^y = \bar{L}, \quad
K^x + K^y = \bar{K}
\]

---

\section{Homogeneous Production Function}

A production function is homogeneous of degree $r$ if:
\[
F(tL, tK) = t^r F(L,K)
\]

If $r=1$, the function exhibits constant returns to scale.

---

\section{Euler's Theorem}

If $F(L,K)$ is homogeneous of degree 1:
\[
F(L,K) = F_L L + F_K K
\]

This implies that total output is exactly exhausted by factor payments.

---

\section{Firm Optimization}

Profit maximization:
\[
\max \pi = pF(L,K) - wL - rK
\]

First-order conditions:
\[
pF_L = w, \quad pF_K = r
\]

Thus:
\[
\frac{F_L}{F_K} = \frac{w}{r}
\]

---

\section{Cost Minimization}

\[
\min C = wL + rK
\quad \text{s.t. } Q = F(L,K)
\]

FOC:
\[
\frac{F_L}{F_K} = \frac{w}{r}
\]

---

\section{Marginal Rate of Technical Substitution}

\[
MRTS = \frac{F_L}{F_K}
\]

At optimum:
\[
MRTS = \frac{w}{r}
\]

---

\section{Capital-Labor Ratio Representation}

Define:
\[
k = \frac{K}{L}
\]

Then:
\[
Q = L f(k)
\]

Marginal products:
\[
F_K = f'(k), \quad
F_L = f(k) - kf'(k)
\]

---

\section{Elasticity of Substitution}

\[
\sigma = \frac{d(K/L)/(K/L)}{d(MRTS)/(MRTS)}
\]

Interpretation:
\begin{itemize}
    \item $\sigma > 1$: factors are easily substitutable
    \item $\sigma < 1$: factors are weak substitutes
\end{itemize}

---

\section{Factor Income Distribution}

From Euler's theorem:
\[
pF(L,K) = wL + rK
\]

Factor shares:
\[
s_L = \frac{wL}{Q}, \quad s_K = \frac{rK}{Q}
\]

---

\section{Production Efficiency}

Efficiency requires:
\[
MRTS^x = MRTS^y = \frac{w}{r}
\]

---

\section{Production Possibility Frontier}

The PPF represents feasible combinations of $(x,y)$.

Slope:
\[
\frac{dy}{dx} = MRT
\]

---

\section{General Equilibrium Conditions}

\subsection{Production Side}
\[
MRTS^x = MRTS^y = \frac{w}{r}
\]

\subsection{Consumption Side}
\[
MRS = \frac{p_X}{p_Y}
\]

\subsection{Overall Condition}
\[
MRS = MRT
\]

---

\section{Price-Factor Relationship}

Unit cost conditions:
\[
p_x = w a_{Lx} + r a_{Kx}
\]

\[
p_y = w a_{Ly} + r a_{Ky}
\]

---

\section{Comparative Statics}

Differentiation:
\[
a_{Lx} dw + a_{Kx} dr = dp_x
\]

\[
a_{Ly} dw + a_{Ky} dr = 0
\]

Matrix form:
\[
\begin{bmatrix}
a_{Lx} & a_{Kx} \\
a_{Ly} & a_{Ky}
\end{bmatrix}
\begin{bmatrix}
dw \\
dr
\end{bmatrix}
=
\begin{bmatrix}
dp_x \\
0
\end{bmatrix}
\]

---

\section{Stolper--Samuelson Theorem}

Results:
\[
\frac{dw}{dp_x} > 0, \quad
\frac{dr}{dp_x} < 0
\]

Interpretation:
\begin{itemize}
    \item Increase in price of labor-intensive good raises wages
    \item Lowers return to capital
\end{itemize}

Stronger result:
\[
\frac{dw}{w} > \frac{dp_x}{p_x}
\]

---

\section{Rybczynski Theorem}

If factor endowment changes:

\[
L = a_{Lx} x + a_{Ly} y
\]
\[
K = a_{Kx} x + a_{Ky} y
\]

Differentiation:
\[
a_{Lx} dx + a_{Ly} dy = d\bar{L}
\]
\[
a_{Kx} dx + a_{Ky} dy = 0
\]

Results:
\[
\frac{dx}{d\bar{L}} > 0, \quad
\frac{dy}{d\bar{L}} < 0
\]

Interpretation:
\begin{itemize}
    \item Increase in labor increases labor-intensive good
    \item Decreases capital-intensive good
\end{itemize}

---

\section{Economic Interpretation}

General equilibrium with production involves:
\begin{itemize}
    \item Resource allocation across industries
    \item Factor price determination
    \item Interaction between goods and factor markets
\end{itemize}